\documentclass[11pt]{article}

\usepackage[margin=1in]{geometry}
\usepackage{amsmath, amssymb, amsthm}
\usepackage{listings}
\usepackage{xcolor}
\usepackage{hyperref}

\hypersetup{colorlinks=true, linkcolor=blue, urlcolor=blue}

% ---- Theorem environments ----
\newtheorem{theorem}{Theorem}
\newtheorem{definition}{Definition}

% ---- Equation numbering to match the paper (19)-(21) ----


% ---- Lean listing style (used only in the final code appendix) ----
\definecolor{leancomment}{RGB}{106,153,85}
\definecolor{leankeyword}{RGB}{150,60,170}
\definecolor{leanbg}{RGB}{246,246,246}

\lstdefinelanguage{Lean}{
  morekeywords={theorem,def,by,apply,intro,exact,rw,simp,fun,trace_state,
    open,import,let,have,show,Type,Prop},
  sensitive=true,
  morecomment=[l]{--},
  morecomment=[s]{/-}{-/},
  morestring=[b]",
}

\lstset{
  language=Lean,
  basicstyle=\ttfamily\small,
  commentstyle=\color{leancomment}\itshape,
  keywordstyle=\color{leankeyword}\bfseries,
  backgroundcolor=\color{leanbg},
  frame=single,
  framesep=6pt,
  breaklines=true,
  showstringspaces=false,
  columns=fullflexible,
  keepspaces=true,
}

% Inline monospace that does NOT trigger listings coloring
\newcommand{\code}[1]{\texttt{#1}}

\newcommand{\R}{\mathbb{R}}
\DeclareMathOperator*{\supp}{sup'}
\DeclareMathOperator*{\infp}{inf'}

\title{Negation Flips \texorpdfstring{$\max$}{max} to \texorpdfstring{$\min$}{min}:\\
  The (20)$\to$(21) Step of Fard \& Pineau (2011), Formalized in Lean}
\author{}
\date{}

\begin{document}
\maketitle

\begin{abstract}
This note isolates and formalizes the single genuine inference in the proof of
Theorem~1 of Fard \& Pineau (2011), \emph{Non-Deterministic Policies in Markovian
Decision Processes} (JAIR 40, pp.~1--24): the move from equation~(20) to
equation~(21), in which negating a maximization is rewritten as a minimization
of the negation. Section~\ref{sec:paper} reproduces the mathematical statement
and its short paper-style proof. Section~\ref{sec:lean} translates the Lean
formalization step by step into ordinary mathematics. The complete Lean~4
(Mathlib) source is collected in the appendix.
\end{abstract}

% ==========================================================================
\section{The Original Result}
\label{sec:paper}

The paper evaluates a non-deterministic policy $\Pi$ by constructing an
\emph{evaluation MDP} $M' = (S, A', R', T, \gamma)$ with $A' = \Pi$ and
$R' = -R$, and proving:

\begin{theorem}[Fard \& Pineau, Thm.~1]
The negated value of the non-deterministic policy $\Pi$ equals the value of the
optimal policy on the evaluation MDP:
\[
  Q^{\Pi}_M(s,a) \;=\; -\,Q^{*}_{M'}(s,a).
\]
\end{theorem}

The proof negates the Bellman optimality equation on $M'$ and pushes the
negation inside the optimization. Written out, the key steps are:
\begin{align}
  Q^{*}_{M'}(s,a)
    &= \bar R'(s,a) + \gamma \sum_{s'} T(s,a,s') \max_{a' \in A'} Q^{*}_{M'}(s',a')
      \tag{19}\\[2pt]
  \Rightarrow\; -Q^{*}_{M'}(s,a)
    &= -\bar R'(s,a) - \gamma \sum_{s'} T(s,a,s') \max_{a' \in A'} Q^{*}_{M'}(s',a')
      \tag{20}\\[2pt]
  \Rightarrow\; -Q^{*}_{M'}(s,a)
    &= \bar R(s,a) + \gamma \sum_{s'} T(s,a,s') \min_{a' \in \Pi(s')}\!\bigl(-Q^{*}_{M'}(s',a')\bigr).
      \tag{21}
\end{align}

Two of the three changes between~(20) and~(21) are purely definitional
substitutions from the construction of $M'$: $-\bar R'(s,a) = \bar R(s,a)$
(since $R' = -R$) and $A' = \Pi$. The one substantive inference is the
transformation
\[
  -\max_{a'} Q^{*}_{M'}(s',a')
  \;=\;
  \min_{a'}\bigl(-Q^{*}_{M'}(s',a')\bigr),
\]
i.e.\ that negation reverses the order and thereby turns the maximizer into the
minimizer. This identity is what Section~\ref{sec:lean} formalizes in full.

Note also a hypothesis the prose leaves implicit but the formalization cannot:
by Definition~1 of the paper, $\Pi(s)$ must be a \emph{non-empty} set of actions.
The quantities $\max$ and $\min$ over $\Pi(s')$ are only well-defined under this
assumption; in Lean it must be supplied explicitly as an argument.

% ==========================================================================
\section{The Formalization, Step by Step}
\label{sec:lean}

We strip the surrounding MDP away and formalize the bare identity. The action
set becomes a finite set $A$ (Mathlib's \code{Finset}), the optimal $Q$-value
$Q^{*}_{M'}(s',\cdot)$ becomes an arbitrary function $Q : A \to \R$, and
$\max/\min$ become Mathlib's $\supp$ and $\infp$ (the primed variants are the
max/min taken over a \emph{non-empty} finite set; the non-emptiness proof
$h$ is a required argument, exactly the Definition~1 hypothesis above). The
claim is
\begin{equation}
  -\,\supp_{a \in A} Q(a) \;=\; \infp_{a \in A}\bigl(-Q(a)\bigr).
  \tag{$\star$}\label{eq:star}
\end{equation}

\subsection*{Overall strategy}

To prove an equality of two reals, prove $\le$ in both directions
(antisymmetry of $\le$). The two directions are mirror images across the
$\max \leftrightarrow \min$ axis. Within each, only two kinds of step occur:
\begin{description}
  \item[Reshape.] A reversible sign-flip identity
    ($-x \le -y \iff y \le x$, or $x \le -y \iff y \le -x$) moves negations
    across $\le$ until the buried $\supp$ or $\infp$ is exposed on the side the
    next fact needs. Truth is unchanged; only the shape of the goal changes.
  \item[Close.] A defining fact about $\max/\min$ discharges a branch: any
    element's value is $\le$ the max; the min is $\le$ any element's value; and
    dually, to bound a min from below (or a max from above) it suffices to bound
    every element.
\end{description}

\subsection*{Direction 1: \texorpdfstring{$-\supp_a Q \le \infp_a(-Q)$}{-sup Q <= inf(-Q)}}

\begin{enumerate}
  \item \textbf{Reduce to elements.} The right-hand side is a minimum. To show
    $c \le \min$, it suffices to show $c \le$ every term the minimum ranges over.
    So the goal becomes: for every $a \in A$,
    \[
      -\,\supp_{a'} Q(a') \;\le\; -Q(a).
    \]
  \item \textbf{Fix an arbitrary $a$.} We prove the statement for a generic
    $a \in A$, using nothing specific about which element it is; this discharges
    all elements at once.
  \item \textbf{Reshape.} Both sides are negated, i.e.\ the goal has the form
    $-x \le -y$. Applying $-x \le -y \iff y \le x$ strips the negations and
    flips the sides:
    \[
      Q(a) \;\le\; \supp_{a'} Q(a').
    \]
  \item \textbf{Close.} This is exactly the defining property of the supremum:
    any element's value is at most the maximum. Done.
\end{enumerate}

\subsection*{Direction 2: \texorpdfstring{$\infp_a(-Q) \le -\supp_a Q$}{inf(-Q) <= -sup Q}}

This is the mirror of Direction~1, with $\max$ and $\min$ swapped.

\begin{enumerate}
  \item \textbf{Reshape \#1.} The goal has the form $x \le -y$ (only the right
    side negated). Applying $x \le -y \iff y \le -x$ moves the negation across,
    exposing the maximum on the left:
    \[
      \supp_{a'} Q(a') \;\le\; -\,\infp_{a'}\bigl(-Q(a')\bigr).
    \]
  \item \textbf{Reduce to elements.} The left-hand side is now a maximum. To
    show $\max \le c$, it suffices to show every term is $\le c$. The goal
    becomes: for every $a \in A$,
    \[
      Q(a) \;\le\; -\,\infp_{a'}\bigl(-Q(a')\bigr).
    \]
  \item \textbf{Fix an arbitrary $a$}, as before.
  \item \textbf{Reshape \#2.} The goal again has the form $x \le -y$. The same
    sign-flip identity now exposes the minimum on the left:
    \[
      \infp_{a'}\bigl(-Q(a')\bigr) \;\le\; -Q(a).
    \]
  \item \textbf{Close.} This is exactly the defining property of the infimum:
    the minimum is at most any element's value --- here the element $-Q(a)$.
    Done.
\end{enumerate}

Both directions established, antisymmetry gives the equality~\eqref{eq:star},
which is precisely the $-\max = \min(-\,\cdot)$ identity underlying the paper's
step (20)$\to$(21). \qed

\subsection*{Two remarks}

\begin{itemize}
  \item \textbf{``Fix an arbitrary $a$'' is not a loop.} A single argument about
    a generic element, using no fact peculiar to it, settles all elements
    simultaneously --- even were $A$ infinite. This is the content of a
    universal proof, and is what the mechanized \code{intro} step performs.
  \item \textbf{Non-emptiness is load-bearing.} Every $\supp$/$\infp$ above
    carries the hidden hypothesis that $A \neq \emptyset$; without it the
    max/min --- and hence the whole identity --- is undefined. The mechanization
    makes this the paper's Definition~1 condition, promoted from prose to an
    unskippable argument.
\end{itemize}

% ==========================================================================
\appendix
\section{Lean 4 (Mathlib) Source}
\label{sec:code}

The complete machine-checked proof. Each comment states the goal \emph{entering}
the line below it; the two \code{apply}/\code{exact} closing lemmas
(\code{le\_sup'}, \code{inf'\_le}) and their bound-from-the-other-side duals
(\code{le\_inf'}, \code{sup'\_le}) are the ``Close'' steps, while
\code{neg\_le\_neg\_iff} and \code{le\_neg} are the ``Reshape'' steps. The proof
checks against Lean's standard axioms only (\code{propext},
\code{Classical.choice}, \code{Quot.sound}).

\begin{lstlisting}
import Mathlib
open Finset

theorem neg_max_eq_min_neg
    {Action : Type*}
    (actions : Finset Action)
    (h_nonempty : actions.Nonempty)
    (Q : Action -> R) :
    - actions.sup' h_nonempty Q
      = actions.inf' h_nonempty (fun x => - Q x) := by
  -- Reduce the equality to two inequalities (<= each direction).
  apply le_antisymm

  -- ===== DIRECTION 1:  - sup' Q  <=  inf' (-Q) =====
  -- Goal: - sup' Q <= inf' (fun x => -Q x)
  -- RHS is a MIN; to bound (something) <= a min, bound it <= every element.
  . apply Finset.le_inf'
    -- Goal: forall b in actions, - sup' Q <= - Q b
    -- Prove the "forall" for an ARBITRARY element (intro destructures the goal:
    -- a : Action, typed from the forall; ha : a in actions).
    intro a ha
    -- Goal: - sup' Q <= - Q a
    -- RESHAPE. Shape is  -x <= -y ; strip both negations and flip sides.
    rw [neg_le_neg_iff]
    -- Goal: Q a <= sup' Q
    -- CLOSE. Exactly "any element's value <= the max".
    exact Finset.le_sup' Q ha

  -- ===== DIRECTION 2:  inf' (-Q)  <=  - sup' Q   (mirror of Direction 1) =====
  -- Goal: inf' (fun x => -Q x) <= - sup' Q
  -- RESHAPE #1. Shape is  x <= -y ; move negation across to expose sup'.
  . rw [le_neg]
    -- Goal: sup' Q <= - inf' (fun x => -Q x)
    -- LHS is a MAX; to bound a max <= (something), bound every element.
    apply Finset.sup'_le
    -- Goal: forall b in actions, Q b <= - inf' (fun x => -Q x)
    -- Arbitrary element again (fresh a/ha, unrelated to Direction 1's).
    intro a ha
    -- Goal: Q a <= - inf' (fun x => -Q x)
    -- RESHAPE #2. Same lemma le_neg; this time it exposes inf' on the left.
    rw [le_neg]
    -- Goal: inf' (fun x => -Q x) <= - Q a
    -- CLOSE. Exactly "the min <= any element's value". Hand inf'_le the function
    -- whose min we mean, plus the membership proof ha. The applied lambda
    -- (fun x => -Q x) a beta-reduces to - Q a, matching the goal.
    exact Finset.inf'_le (fun x => - Q x) ha
\end{lstlisting}

\end{document}
